%!TEX root=../a03.tex
\section{Q5}
\begin{enumerate}[label = (\alph*)]
    \item \begin{solution}
        In order to prove that $K \cap L \neq \emptyset$. We need to show that there exists $x$ such that 
        $x \in K$ and $x \in L$. By the definition of sets $K$ and $L$, it suffices to prove that 
        \begin{equation*}
            \exists x, s, u \in \Z, 3000 \leq s \leq 4000, x = \paren{2s}^2 = 2u^3. 
        \end{equation*} 

        We pick $s = 2 \cdot \paren{12}^3$ and $u = 2 \cdot 12^2$

        First, we need to prove that $3000 \leq s \leq 4000$. 
        \begin{equation*}
            s = 2 \cdot 1728 = 3456,  
        \end{equation*}

        which satisfies $3000 \leq 3456 \leq 4000$. 

        Next, we need to prove $\paren{2s}^2 = 2u^3$. Subtituting in $s = 2 \cdot \paren{12}^3$ and $u = 2 \cdot 12^2$, we have 
        \begin{equation*}
            \paren{2s}^2 = \bracks{2 \cdot \paren{2 \cdot 12^3}}^2 = 2^4 \cdot 12^6  = 2 \paren{2^3 \cdot 12^6}= 2 \paren{2 \cdot 12^2}^3 = 2u^3. 
        \end{equation*}

        Since $x=(2s)^2$ with $s\in\Z$ and $3000\leq s\leq 4000$, we have $x\in K$; moreover, since $x=2u^3$ with $u\in\Z$, we have $x\in L$.

        Thus, 
        \begin{equation*}
            x \in \paren{K \cap L}
        \end{equation*}

        Hence,  we prove that sets $K$ and $L$ are not disjoint. 
    \end{solution}

    \newpage

    \item \begin{solution}
        The following statement is true.

In order to prove the statement, by definition of the sets, we need to show that there is exactly one pair of integers $s,u$ such that
\begin{equation*}
    3000 \leq s \leq 4000,
    (2s)^2 = 2u^3.
\end{equation*}

Assume $(2s)^2 = 2u^3$.
We have
\begin{align*}
    (2s)^2 = 2u^3
    &\implies 4s^2 = 2u^3 \\
    &\implies 2s^2 = u^3.
\end{align*}

According to the lemma 2 shown after the solution for Question (b), there exists
$k \in \Z$ such that
\begin{equation*}
    u = 2k.
\end{equation*}

Substituting this back into $2s^2 = u^3$, we get
\begin{align*}
    2s^2 = (2k)^3
    &\implies 2s^2 = 8k^3 \\
    &\implies \frac{1}{4}s^2 = k^3 \\
    &\implies \left(\frac{1}{2}s\right)^2 = k^3.
\end{align*}

Since $s^2 = 2\paren{2k^3}$, by the lemma 1, $s$ is even, so $\frac{s }{2 } \in \Z$. 
By the proposition given for this question, $k$ is a perfect square. 
Therefore, there exists some non-negative $t \in \Z$ such that
\begin{equation*}
    k = t^2.
\end{equation*}

Substituting $k = t^2$ into the equation, we get
\begin{align*}
    \left(\frac{1}{2}s\right)^2
    &= (t^2)^3 \\
    &= t^6.
\end{align*}

Since $s$ and $t$ are both non-negative, taking the square root of both sides gives
\begin{align*}
    \frac{1}{2}s &= t^3 \\
    \implies s &= 2t^3.
\end{align*}

In order for $(2s)^2 \in K$, we must have
\begin{equation*}
    3000 \leq s \leq 4000.
\end{equation*}

Subtituting $s=2t^3$ in,
\begin{equation*}
    3000 \leq 2t^3 \leq 4000,
\end{equation*}
\begin{equation*}
    \implies 1500 \leq t^3 \leq 2000.
\end{equation*}

Try the first case,
\begin{equation*}
    11^3 = 1331 < 1500
\end{equation*}
and the second case
\begin{equation*}
    13^3 = 2197 > 2000.
\end{equation*}

There is only one $t \in \N$ satisfies 
\begin{equation*}
    1500 \leq t^3 \leq 2000, 
\end{equation*}
which is
\begin{equation*}
    t = 12.
\end{equation*}

Therefore,
\begin{equation*}
    s = 2(12)^3 = 3456
\end{equation*}
and
\begin{equation*}
    u = 2k = 2t^2 = 2(12)^2 = 288.
\end{equation*}

Thus, the only element in both sets $K$ and $L$ is
\begin{equation*}
    x = (2s)^2 = (2\cdot 3456)^2 = 2(288)^3.
\end{equation*}

it follows that
\begin{equation*}
    |K \cap L| = 1.
\end{equation*}

Therefore, the statement is true.
    \end{solution}
\end{enumerate}

\begin{lemma}
    \begin{equation*}
        \forall n \in \Z, \bracks{\paren{2 \mid n^2} \implies \paren{2 \mid n}}.
    \end{equation*}
\end{lemma}
\begin{proof}
    Suppose n is arbitrary in $\Z$, for contrapositive, $2 \nmid n$. 

    By the definition of divisibility, there exists $k \in \Z$ such that 
    \begin{equation*}
        n = 2k + 1. 
    \end{equation*}

    This implies 
    \begin{equation*}
        n^2 = \paren{2k + 1}^2 = 4k^2 + 4k + 1
    \end{equation*}

    Since $2k^2 + 2k \in \Z$, we have 
    \begin{equation*}
        n^2 = 2(2k^2 + 2k) + 1
    \end{equation*}

    By the definition of divisibility, we get 
    \begin{equation*}
        2 \nmid n^2. 
    \end{equation*}

    Thus, we prove that if $2 \mid n^2$, then $2 \mid n$. 
\end{proof}

\begin{lemma}
    \begin{equation*}
        \forall n \in \Z, \bracks{\paren{2 \mid n^3} \implies \paren{2 \mid n}}.
    \end{equation*}
\end{lemma}

\begin{proof}
    Suppose $n$ is arbitray in $\Z$, for contrapositive, $2 \nmid n$. 

    By the definition of divisibility, there exists $k \in \Z$ such that 
    \begin{equation*}
        n = 2k + 1. 
    \end{equation*}

    This implies 
    \begin{equation*}
        n^3 = \paren{2k + 1}^3 = 8k^3 + 12k^2 + 6k + 1.
    \end{equation*}

    Since $4k^3 + 6k^2 + 3k \in \Z$, we have 
    \begin{equation*}
        n^3 = 2(4k^3 + 6k^2 + 3k) + 1.
    \end{equation*}

    By the definition of divisibility, we get 
    \begin{equation*}
        2 \nmid n^3. 
    \end{equation*}

    Thus, we prove that if $2 \mid n^3$, then $2 \mid n$. 
\end{proof}